\section{Discussion and outlook}
\label{sec:discussion}

We propose Murmur to connect agent communication with organisational commitment, and hypothesise that its shared interaction records improve independent reconstruction and recovery. A conversation can become a common source of evidence when its participants, dependencies and outcomes are published in a form that authorised parties can retain and verify. Whether that representation improves coordination over well-designed alternatives remains open. The proposal is useful insofar as it makes the required comparison and its failure conditions precise.

The strongest competing explanation is that the familiar enforcement mechanisms do the consequential work. Atomic reservations prevent overspending; unique terminal decisions prevent contradictory local outcomes; durable outboxes preserve decisions across publication failures. These properties can be implemented with other transports. The finite-model results are consistent with Murmur's design, but do not discriminate between these explanations. A demonstrated benefit must therefore concern evidence discovery, reconstruction or recovery after those mechanisms are held equal.

The same records that support reconstruction create a disclosure obligation. Replication can leave a released message beyond the author's control, while restricted access may prevent another participant from seeing evidence needed to interpret a decision. The proposed release interface narrows this tension for approved work orders. It does not establish a general solution for private, adaptive negotiation. A useful coordination substrate must be evaluated both for the evidence it makes available and for the information it unnecessarily exposes.

The likely value of the approach depends on deployment. Stable workflows within one organisation may not justify federation. Changing participant sets and independently operated services provide a stronger setting for testing the hypothesis, because evidence must survive beyond the originating agent and its local infrastructure. An optional runtime such as Agent Substrate changes execution and isolation, but cannot by itself establish this coordination benefit. Its effect should be evaluated separately from the record representation.

\section{Conclusion}
\label{sec:conclusion}

We propose Murmur as a coordination method for agent-mediated business. It separates published coordination from A2A task execution and connects exact commercial terms to delegated authority, resource reservations and controlled release of information. We hypothesise that its shared, attributable interaction records improve how independently operated agents reconstruct and recover business decisions. These controls complement the record: attributable messages alone cannot prevent overspending or disclosure.

The sandbox demonstrates the coordination path, while the worked transaction and bounded model checks examine the proposed commercial rules. They do not establish an advantage across independent organisations. That claim requires comparison with alternatives providing equivalent authority, resource and retention controls. The central question is whether shared records improve reconstruction and recovery without unacceptable disclosure or operational cost. Answering it would establish when a conversation should become part of the infrastructure for agent-mediated business.
